\documentclass[11pt]{article}
\usepackage[T1]{fontenc}
\usepackage{amsmath}
\usepackage{amssymb}
\usepackage{amsthm}
\usepackage[hidelinks]{hyperref}
\newtheorem{theorem}{Theorem}[section]
\newtheorem{lemma}[theorem]{Lemma}
\newtheorem{corollary}[theorem]{Corollary}
\theoremstyle{definition}
\newtheorem{definition}[theorem]{Definition}
\begin{document}
\begin{center}
{\LARGE Natural number addition}
\end{center}
\section{Additive identities}
\label{ll:addition:identities}
\begin{theorem}
\label{ll:addition:add-zero}
For every natural number \(n\), \(n + 0 = n\).
\end{theorem}
\begin{proof}
The goal follows by reflexivity.
\end{proof}
\begin{theorem}
\label{ll:addition:zero-add}
For every natural number \(n\), \(0 + n = n\).
\end{theorem}
\begin{proof}
The goal follows from \(\operatorname{zeroadd}(n)\).
\end{proof}
\begin{lemma}
\label{ll:addition:succ-add}
For every natural number \(n\) and natural number \(m\), \(\operatorname{succ}(n) + m = \operatorname{succ}(n + m)\).
\end{lemma}
\begin{proof}
The goal follows from \(\operatorname{succadd}(n, m)\).
\end{proof}
\begin{corollary}
\label{ll:addition:both-identities}
For every natural number \(n\), \(n + 0 = n\) and \(0 + n = n\).
\end{corollary}
\begin{proof}
Branch \(1\):
\begin{quote}
The goal follows by reflexivity.
\end{quote}
Branch \(2\):
\begin{quote}
The goal follows from \(\texttt{Addition::zero-add}(n)\).
\end{quote}
\end{proof}
\section{Commutativity}
\label{ll:addition:commutativity}
\[\mathrm{Parameters}: \forall a \in \mathbb{N}; \forall b \in \mathbb{N}\]
\begin{theorem}
\label{ll:addition:add-comm}
\(a + b = b + a\).
\end{theorem}
\begin{proof}
The goal follows from \(\operatorname{addcomm}(a, b)\).
\end{proof}
\begin{corollary}
\label{ll:addition:add-comm-symm}
\(b + a = a + b\).
\end{corollary}
\begin{proof}
The goal follows from \(\operatorname{eqsymm}(\texttt{Addition::add-comm}(a, b))\).
\end{proof}
\section{Associativity}
\label{ll:addition:associativity}
\begin{theorem}
\label{ll:addition:add-assoc}
For every natural number \(a\) and natural number \(b\) and natural number \(c\), \(a + b + c = a + (b + c)\).
\end{theorem}
\begin{proof}
The goal follows from \(\operatorname{addassoc}(a, b, c)\).
\end{proof}
\begin{theorem}
\label{ll:addition:add-left-comm}
For every natural number \(a\) and natural number \(b\) and natural number \(c\), \(a + (b + c) = b + (a + c)\).
\end{theorem}
\begin{proof}
Rewrite the goal using \(\leftarrow \texttt{Addition::add-assoc}(a, b, c)\), \(\rightarrow \texttt{Addition::add-comm}(a, b)\), \(\rightarrow \texttt{Addition::add-assoc}(b, a, c)\).
\end{proof}
\end{document}
